\begin{abstract}
When does running a block of layers again work as test-time compute? A looped transformer adds depth without adding
parameters, and two benefits are claimed for it: that it behaves like a larger model, and that more loops at inference
solve harder problems without a written chain of thought. We test both: on web text in a pre-registered study, and on a
synthetic serial task whose design and predictions were frozen before its main runs. On web text we
trained a grid of 111 models from 10M to 80M parameters with dense references up to 160M, varying where the loop sits,
how often it runs and how it is backpropagated, and extended it in model size, data and learning rate. On a
lookup-chain task that needs serial steps we compared looped models with dense models of equal parameters and of
equal executed depth, with a written chain of thought and with pause tokens. None of the four pre-registered
hypotheses held as stated, and the pre-registered effective-parameter exponent is not identified. In an exploratory
analysis without a fitted law, each parameter of a model whose middle block is looped 4 times is worth at most about
one dense parameter at 10--20M with up to 20 training tokens per parameter and about 1.3--1.5 at 40M--160M with 20--40,
rising with the tokens per parameter; at equal training compute and a shared learning rate the dense model wins every measured comparison, by a
margin that shrinks with data. On text the
looped block runs a fixed-length program, a cycle or a fixed point, depending on training, and gains nothing from
extra loops. On the chain task the same architecture computes the chain step by step inside the loop and solves
longer chains when run longer (a model trained with 8 loops on chains of at most 8 steps goes from 0.04 to 0.83
accuracy on 12-step chains when run with 16 loops), and for most training regimes its own confidence tells it when to
stop. Dense models of
equal depth do not extend beyond the training range, and neither does a chain of thought on this task, whose prompt
grows with the chain. At deployment and equal loss, a middle block looped 4 times saves at most 19\% of the weight
memory and costs about 2.2--2.5 times the KV cache and the decoding time. Code, configurations, training logs and results are released.
\end{abstract}
